
简单来说，深度学习就是一种包括多个隐含层（越多即为越深）的多层感知机。它通过组合低层特征，形成更为抽象的高层表示，用以描述被识别对象的高级属性类别或特征。能自生成数据的中间表示（虽然这个表示并不能被人类理解），是深度学习区别于其他机器学习算法的独门绝技。


%%%%%  以下可能不需要
深度学习与强化学习：机器学习的新纪元
深度学习的出现
深度学习通过多层神经网络对数据进行复杂特征的分层提取，显著提高了模型对高维复杂数据的处理能力。

它使得无监督学习在特征学习和生成建模中获得了新的突破。
它在监督学习中实现了端到端的训练，大幅提升了传统任务的性能。

强化学习的崛起
强化学习通过试探和反馈，解决了传统学习方法无法应对的动态决策问题，例如游戏 AI（AlphaGo）、自动驾驶和机器人控制。

机器学习的新范式
深度学习和强化学习的结合，使得机器学习从“标注依赖的模型”转变为“主动学习的智能体”，能够像生物大脑一样从环境中自我学习并适应。
